\documentclass[12pt]{article}

\usepackage[a4paper, total={6in, 8in}]{geometry}
\usepackage{textcomp}

\begin{document}
\setlength{\parindent}{0pt}

\def \t {\textrightarrow}
\def \v {\vspace{1em}}
\def \bi {\begin{itemize}}
\def \ei {\end{itemize}}
\def \s[#1] {\section*{#1}}
\def \ss[#1] {\subsection*{#1}}
\def \sss[#1] {\subsubsection*{#1}}

\s[D'annunzio]
Dotato di sensibilita particolare, indirizzata in abitudini particolari \\
Stile di vita che si riassume in un detto paradigmatico: "vivere la vita come un opera d arte", il vivere inimitabile \\
Forte connessione tra vita e arte \\
Voyerismo: era vorace della vita, e viaggiava con la mente \t con tendenza a vivere ogni esperienza intensamente, e la tendenza a essere una figura mediatica e out of line, sempre sotto la luce dei riflettori \\
Leopardi facile all'esposizione della sua anima, pascoli del suo vissuto, d annunzio alla sua vita 

\ss[Vita]
Nasce da famiglia borghese nel 1863 \t e ha formazione scolastica rigida, frequenta Coleggio Cicognini a Prato (istituto + autorevole del tempo) \t ma è ribelle \\
A gia 16 anni pubblica la prima raccolta, nel 79, chiamato "Primo vere" ? \\
Successo in realta anche grazie al fatto che aveva fatto scrivere a un giornale  della sua morte, fa notizia e quindi si diffonde la sua opera \\
Studia poi all universita di roma \t ma abbandona presto i progetti di laurea \t privilegia i salotti della roma del tempo \\
Non è la roma di Pasolini, viscerale e genuina, ma elitaria? \\
Roma è la citta della classicita e del rinnovamento \t e gli anni di roma sono gli anni che mostrano la sua propensione per lo scandalo \\
Vive vita mondana senza freni, di lusso (come l uomo di lusso di verga) \t questo porta all'indebitarsi, e quindi fugge a napoli (dove conduce vita infuocata) \\
Prende forma della sua identita di esteta \t lui si esibisce \\
Anche da Napoli dovra scappare perche indebitato
\end{document}
